% ---------------------------------------------------------------------------
% howtoread.tex
% ---------------------------------------------------------------------------
\chapter*{How to Read This Booklet}
\addcontentsline{toc}{chapter}{How to Read This Booklet}
\markboth{How to Read This Booklet}{How to Read This Booklet}

\section*{Typographic conventions}

\begin{itemize}
  \item \code{Monospaced text} is code: an identifier, a keyword, a file name
        or a command.
  \item \kw{Coloured monospaced text} marks a language keyword under
        discussion.
  \item \file{Green monospaced text} is a file in the accompanying source
        tree.
  \item \tool{Sans-serif} names are tools: \tool{QuestaSim}, \tool{Verilator},
        \tool{cocotb}, \tool{GTKWave}.
\end{itemize}

\section*{The three kinds of box}

\begin{gotcha}{What a trap looks like}
A trap marks a mistake that VHDL experience makes you \emph{more} likely to
commit, not less. They are the part worth re-reading before your first
design review.
\end{gotcha}

\begin{tipbox}{What a recipe looks like}
A recipe is a pattern to copy. When the text says ``write it like this'',
this is the box it is in.
\end{tipbox}

\begin{notebox}{What a note looks like}
A note supplies background: where a construct came from, what the standard
actually says, why a tool behaves the way it does. Skip these on a first
reading and come back when something surprises you.
\end{notebox}

\begin{ruleofthumb}
A rule of thumb is one line you should be able to recite from memory.
\end{ruleofthumb}

\section*{If you have two weeks before the course starts}

Read \cref{ch:philosophy} carefully. Skim \cref{ch:rosetta} once, without
trying to memorise it, so that you know what is in it. Then move on to
\refverificationbook for the testbench, and to \refexamplebook to see it all
used together.

\section*{If you are converting an existing VHDL design}

Start at \cref{ch:rosetta} and work through your source file construct by
construct. Pay particular attention to the sections on widths and
signedness, which is where mechanical translation goes wrong. Then read
\cref{ch:beyond} for the constructs that will let you delete half of what
you just translated.

\section*{The companion booklets}

This is the \textbf{Design Cookbook}. \refverificationbook covers everything
that has no VHDL counterpart: classes, randomisation, coverage, assertions,
UVM. \refsimulationbook covers \tool{Verilator}, \tool{QuestaSim} and
Python/\tool{cocotb} testbenches. \refexamplebook works one CORDIC design
through all of it.
